\begin{corollary}[Proved: residual genus and different bounds]
\label{D:surface-residual-genus-bound}
Use the resolution setup of
\ref{D:surface-uniform-different}, or its fixed-boundary variant
\ref{D:fixed-boundary-uniform-different}. Write the smooth exceptional
curves as $E_i$, choose a positive integral exceptional divisor
$A=\sum_i a_iE_i$ with $A\cdot E_i<0$, and put
\[
 C_A=(K_Y+D)\cdot A+\sum_i a_i,
\]
where $D=0$ in the punctured case and $D$ is the reduced strict
transform of the original divisorial boundary in the tame variant.

Let a normal finite covering in the preceding etale/tame setup
(hence generically separable) have generic degree $n$, and form
$f:Z\to Y$ as in those results, contracting only rational
$(-1)$-curves mapped to points of $Y$. For every component $F$
dominating $E_i$, let $e_F$ be its ramification index and let
$g_F$ be the genus of its normalization. Then
\[
 2\sum_{F/E_i}a_i e_Fg_F\leq n C_A.
 \tag{15}
\]
If $Q$ ranges over the remaining point-mapped exceptional components,
$m_Q$ is its coefficient in $f^*A$, and $g_Q$ is its normalized
genus, the stronger inequality
\[
 2\sum_{F/E_i}a_i e_Fg_F+\sum_Qm_Qg_Q\leq n C_A
 \tag{16}
\]
holds.

Factor the finite residual curve map from the normalization of $F$
to $E_i$ into a purely inseparable map followed by a separable map
$C_F\to E_i$, and let $\Delta_F$ be the different divisor of the
separable map. Then
\[
 \sum_{F/E_i}a_i e_F\deg\Delta_F
 \leq n\left(C_A+2\sum_{g(E_i)=0}a_i\right).
 \tag{17}
\]
All constants depend only on the fixed resolution and boundary.
\end{corollary}

\begin{proof}
Every component of the exceptional fibre has negative integral
self-intersection. Adjunction therefore gives, for a dominating
component,
\[
 K_Z\cdot F=2p_a(F)-2-F^2\geq2g_F-1.
\]
For a point-mapped component $Q$, the contractions stipulated in
the setup ensure $K_Z\cdot Q\geq0$. If $g_Q\geq1$, adjunction
also gives $K_Z\cdot Q\geq2g_Q-1\geq g_Q$; hence this last
inequality holds for every such $Q$.

The coefficient of $f^*A$ along $F/E_i$ is $a_i e_F$.
The Jacobian and projection calculation in the cited results gives
\[
 K_Z\cdot f^*A\leq n(K_Y+D)\cdot A.
\]
Here the vertical Jacobian contribution has nonpositive
intersection with $f^*A$; the tame horizontal contribution is
bounded by $nD\cdot A$. Thus
\[
 2\sum a_i e_Fg_F+\sum_Qm_Qg_Q
 \leq n(K_Y+D)\cdot A+\sum a_i e_F.
\]
At the generic point of $E_i$,
$\sum_{F/E_i}e_F[k(F):k(E_i)]=n$, so
$\sum_{F/E_i}e_F\leq n$. This proves (16), and hence (15).

The ground field is algebraically closed. Purely inseparable maps
between smooth proper curves are iterated relative Frobenius maps,
so $g(C_F)=g_F$; see
\citep[Lemmas~53.13.4--53.13.5 and Proposition~53.13.7]{StacksResidualGenusB}.
If $s_F$ is the separable residual degree,
Riemann--Hurwitz \citep[Lemma~53.12.2]{StacksResidualGenusB} gives
\[
 \deg\Delta_F=2g_F-2-s_F(2g(E_i)-2).
\]
For $g(E_i)\geq1$ the last term is nonpositive. For a rational
$E_i$, its positive contribution is $2s_F$, and
$\sum_{F/E_i}e_Fs_F\leq n$. Multiplying by $a_i$, summing,
and using (15) proves (17).
\end{proof}

\begin{remark}[The remaining counting issue]
These inequalities control the total normalized genus and
separable residual ramification by a constant times the covering
degree. They do not fix the residual branch locus or control all
point-mapped rational configurations: a rational $(-2)$-curve
has zero canonical intersection and contributes nothing to
(16). Consequently no bound on the number of coverings, and no
finite-generation conclusion, is inferred from (15)--(17) alone.
\end{remark}
